\documentclass{report}
\usepackage{amsmath,amssymb}
\setlength{\parindent}{0mm}
\setlength{\parskip}{1em}
\begin{document}
\begin{center}
\rule{6in}{1pt} \
{\large Jeffery Allen \\
{\bf New FOSLS Formulation of Nonlinear Stokes Flow for Glaciers}}

2210 Walnut Apt 1 \\ Boulder \\ CO 80302
\\
{\tt jeffery.allen@colorado.edu}\\
Thomas Manteuffel \\
Harihar Rajaram\end{center}

This talk describes two First-order System Least-squares (FOSLS)
formulations of the nonlinear Stokes flow used to model glaciers and ice
sheets. The first is a Stress formulation and the second
a Stress-Vorticity formulation. Both use fluidity, which is the
reciprocal of viscosity and avoid the difficulties of infinite viscosity.
Coercivity and continuity in appropriate Sobolev norms will be discussed.
A Nested Iteration (NI), Newton-FOSLS-AMG approach is employed, in which
the majority of the work
is done on coarse grids.

In current ice sheet codes, viscosity is modeled using a formula known as
Glen's law. It becomes infinite as the gradient of velocity approaches
zero, which can happen near the bed and along the top surface, where ice
is carried along undeformed. In current ice sheet codes, this is
mitigated by a small modification to Glen's law that bounds the maximum
value of viscosity. The fluidity formulation presented here avoids this
difficulty and requires no modification to Glen�s law.

Numerical tests are presented that demonstrate the efficacy of the
NI-Newton-FOSLS-AMG approach on model problems from the ISMIP-HOM problem
set.


\end{document}
